% \documentclass[aspectratio=169,notes]{beamer}
\documentclass[aspectratio=169]{beamer}
\usetheme[faculty=phil]{fibeamer}
\usepackage{polyglossia}
\setmainlanguage{english} %% main locale instead of `english`, you
%% can typeset the presentation in either Czech or Slovak,
%% respectively.
\setotherlanguages{russian} %% The additional keys allow
%%
%%   \begin{otherlanguage}{czech}   ... \end{otherlanguage}
%%   \begin{otherlanguage}{slovak}  ... \end{otherlanguage}
%%
%% These macros specify information about the presentation
\title[Artificial Intelligence Software (AIS)]{Artificial Intelligence Software, Project 1} %% that will be typeset on the
\subtitle{Fine-tuning YOLO models on custom datasets  
\\ \  \\ \  } %% title page.
\author{Oleg Bulichev}
%% These additional packages are used within the document:
\usepackage{ragged2e}  % `\justifying` text
\usepackage{booktabs}  % Tables
\usepackage{tabularx}
\usepackage{tikz}      % Diagrams
\usetikzlibrary{calc, shapes, backgrounds}
\usetikzlibrary{decorations.pathreplacing,calligraphy,calc,graphs}
\usepackage{amsmath, amssymb}
\usepackage{url}       % `\url`s
\usepackage{listings}  % Code listings
% \usepackage{subfigure}
\usepackage{floatrow}
\usepackage{subcaption}
\usepackage{mathtools}
\usepackage{todonotes}
\usepackage{fontspec}
\usepackage{multicol}
\usepackage{pdfpages}
\usepackage{wrapfig}
\usepackage{qrcode}
\usepackage{animate}
\usepackage{booktabs}
\usepackage{multirow}
% \usepackage{graphicx}
\usepackage{colortbl}

\graphicspath{{resources/}}
\frenchspacing

\setbeamertemplate{caption}[numbered]
\usetikzlibrary{graphs}

% \usepackage[backend=biber,style=ieee,autocite=footnote]{biblatex}
% \addbibresource{biblio.bib}
% \DefineBibliographyStrings{english}{%
%   bibliography = {References},}

\newcommand{\oleg}[2][] {\todo[color=red, #1] {OLEG:\\ #2}}
\newcommand{\fbckg}[1]{\usebackgroundtemplate{\includegraphics[width=\paperwidth]{#1}}}%frame background

\usepackage[framemethod=TikZ]{mdframed}
\newcommand{\dbox}[1]{
\begin{mdframed}[roundcorner=3pt, backgroundcolor=yellow, linewidth=0]
\vspace{1mm}
{#1}
\vspace{1mm}
\end{mdframed}
}

\begin{document}
\setlength{\abovedisplayskip}{0pt}
\setlength{\belowdisplayskip}{0pt}
\setlength{\abovedisplayshortskip}{0pt}
\setlength{\belowdisplayshortskip}{0pt}

\fbckg{fibeamer/figs/title_page.png}
\frame[c]{\setcounter{framenumber}{0}
    \usebeamerfont{title}%
    \usebeamercolor[fg]{title}%
    \begin{minipage}[b][6.5\baselineskip][b]{\textwidth}%
        \textcolor{black}{\raggedright\inserttitle}
    \end{minipage}
    % \vskip-1.5\baselineskip

    \usebeamerfont{subtitle}%
    \usebeamercolor[fg]{framesubtitle}%
    \begin{minipage}[b][3\baselineskip][b]{\textwidth}
        \raggedright%
        \insertsubtitle%
    \end{minipage}
    \vskip.25\baselineskip
}
%   \frame[c]{\maketitle}

\fbckg{fibeamer/figs/common.png}

\note{\scriptsize \begin{itemize}
        \item \
    \end{itemize}}

\begin{frame}[t]{Short Description}
The goal of this assignment is fine-tuning a pretrained \textbf{Ultralytics YOLO model} on a custom dataset for a computer vision task of your choice. It's important to implement a \textbf{methodologically correct ML workflow}.
The workflow must include:
\begin{multicols}{2}
    \begin{itemize}
        \item dataset collection
\item dataset annotation
\item dataset versioning
\item experiment tracking
\item model training and improvement
\item deployment and demonstration
    \end{itemize}

You must use:
\begin{itemize}
\item \textbf{Ultralytics YOLO}
\item \textbf{ClearML} for dataset versioning and experiment tracking
\end{itemize}
\end{multicols}

\end{frame}


\begin{frame}[t]{1. Task Selection}
    \vspace*{-0.4cm}
Choose a computer vision task that can be demonstrated using:

\begin{itemize}
\item a \textbf{phone camera}, or
\item a \textbf{webcam on your computer}
\end{itemize}

Students choose the task themselves.

Possible task types:
\begin{itemize}
\item object detection
\item instance segmentation
\item image classification
\end{itemize}

Important requirement:

The object or class must be \textbf{physically accessible}, since you will
collect and annotate your own data.
\end{frame}


\begin{frame}[t]{2. Dataset Preparation}
    \vspace*{-0.4cm}
You must use both:
\begin{multicols}{2}
\begin{itemize}
\item an \textbf{open dataset}
\item a \textbf{self-collected dataset}
\end{itemize}

Possible sources of open datasets:
\begin{itemize}
\item Kaggle
\item Roboflow Universe
\item Hugging Face Datasets
\item other public datasets
\end{itemize}

Your own dataset must contain:

\begin{itemize}
\item \textbf{100--300 images minimum}
\item different lighting conditions
\item different backgrounds
\item different object poses and viewpoints
\end{itemize}
\end{multicols}

\end{frame}


\begin{frame}[t]{3. Data Annotation and Versioning}
All collected images must be \textbf{annotated}.

Possible annotation tools:
\begin{itemize}
\item CVAT
\item Roboflow
\end{itemize}

Annotations must be exported in \textbf{Ultralytics YOLO format}.

Datasets must be \textbf{versioned using ClearML}.

You must create at least:
\begin{itemize}
\item one \textbf{initial dataset version}
\item one \textbf{improved dataset version} (after additional data collection or re-annotation)
\end{itemize}
\end{frame}


\begin{frame}[t]{4. Baseline Model}
Train a \textbf{baseline model} by fine-tuning a pretrained
\textbf{Ultralytics YOLO model}.

Students may choose the model architecture depending on the task.

Examples:
\begin{itemize}
\item smaller models for faster inference
\item larger models if computational resources allow
\end{itemize}

All training runs must be \textbf{logged in ClearML}.
\end{frame}


\begin{frame}[t]{5. Evaluation Metrics}
Use the \href{https://docs.ultralytics.com/metrics/}{standard Ultralytics evaluation metrics}

\begin{itemize}
\item Precision
\item Recall
\item mAP
\item F1 score
\item Confusion Matrix
\end{itemize}

Metrics must be tracked through \textbf{ClearML experiments}.
\end{frame}


\begin{frame}[t]{6. Experiments and Improvements}
You must run at least \textbf{three experiments} aimed at improving the model.
\begin{multicols}{2}
Possible approaches:

\begin{itemize}
\item hyperparameter tuning
\item data augmentation
\item adding more training data
\item fixing annotation errors
\item balancing classes
\item trying different YOLO model sizes
\end{itemize}

Each experiment must be:
\begin{itemize}
\item logged in \textbf{ClearML}
\item reproducible
\end{itemize}
\end{multicols}

\end{frame}


\begin{frame}[t]{7. Training Resources}
This assignment \textbf{can be completed on a CPU-only computer}.

Using additional resources is allowed but not required:

\begin{itemize}
\item GPU
\item Google Colab
\item remote servers
\end{itemize}

Since you are performing \textbf{fine-tuning}, training on CPU is feasible.
\end{frame}


\begin{frame}[t]{8. Deployment and Demonstration}
Export the trained model using \textbf{Ultralytics}.

Demonstrate the solution in real time using:

\begin{itemize}
\item a \textbf{webcam}, or
\item a \textbf{phone camera}
\end{itemize}

The system must show \textbf{live inference} detecting the selected objects.
\end{frame}


\begin{frame}[t]{Submission}
You must submit a \textbf{report} describing the full workflow.

The report must include:
\begin{multicols}{2}
\begin{itemize}
\item description of the selected task
\item links to \textbf{ClearML project and experiments}
\item dataset sources and dataset versions
\item description of experiments
\item evaluation metrics
\item description of improvements
\item demonstration of the final system (video or live demo)
\item discussion of \textbf{difficulties and limitations}
\end{itemize}
\end{multicols}


The goal is to document the \textbf{entire ML workflow}.
\end{frame}




\fbckg{fibeamer/figs/last_page.png}
\frame[plain]{}

\end{document}